\documentclass[border=0pt,varwidth=false]{standalone}
\usepackage{ews-figures}
\begin{document}
\begin{ewsfigure}

\begin{ewsdiagram}
\begin{tikzpicture}[
  every node/.style={font=\fontsize{7.6}{9}\selectfont},
  stage/.style={draw=ewsrule, fill=white, rounded corners=1.5pt,
                inner sep=4pt, minimum height=17pt, minimum width=54pt,
                align=center},
  host/.style={stage, draw=ewsslate, fill=ewsslate!6},
  flow/.style={-{Stealth[length=4.5pt]}, draw=ewsmute, line width=0.6pt},
  pin/.style={draw=ewsaccent, fill=white, rounded corners=1.5pt,
              inner sep=3pt, minimum height=13pt, align=center,
              font=\fontsize{7}{8.4}\selectfont},
  tick/.style={draw=ewsaccent, line width=0.5pt, densely dotted}]

% ---- the deployed pipeline, left to right
\node[stage] (inp) {input\\representation};
\node[stage, right=13pt of inp] (dia) {dialogue\\state};
\node[stage, right=13pt of dia] (rea) {reasoning\\field};
\node[stage, right=13pt of rea] (fin) {final\\answer};
\node[stage, right=13pt of fin] (too) {tool\\call};
\node[host,  right=13pt of too] (hos) {host\\action};

\foreach \a/\b in {inp/dia, dia/rea, rea/fin, fin/too, too/hos}
  \draw[flow] (\a) -- (\b);

% ---- where each finding crosses a boundary
\node[pin, below=22pt of inp] (f1) {\textbf{F1} multilingual\\saturation};
\node[pin, below=22pt of dia] (f2) {\textbf{F2} stateful\\escalation};
\node[pin, below=22pt of rea] (f3) {\textbf{F3} cross-channel\\containment};
\node[pin, below=22pt of too] (f5) {\textbf{F5} code-bearing\\tool call};

\foreach \s/\f in {inp/f1, dia/f2, rea/f3, too/f5}
  \draw[tick] (\s.south) -- (\f.north);

% F4 spans dialogue state and reasoning rather than sitting at one stage.
% Route via a bus above the box so no connector crosses the label.
\coordinate (busl) at ($(f2.south)+(0,-9pt)$);
\coordinate (busr) at ($(f3.south)+(0,-9pt)$);
\coordinate (busm) at ($(busl)!0.5!(busr)$);
\node[pin, anchor=north] (f4) at ($(busm)+(0,-9pt)$)
  {\textbf{F4} persona and agent state across two memory topologies};
\draw[tick] (f2.south) -- (busl);
\draw[tick] (f3.south) -- (busr);
\draw[tick] (busl) -- (busr);
\draw[tick] (busm) -- (f4.north);

% ---- the boundary the paper insists on
\draw[draw=ewsslate, densely dashed, line width=0.5pt]
  ($(too.north east)+(6.5pt,7pt)$) -- ($(too.south east)+(6.5pt,-58pt)$);
\node[anchor=west, font=\fontsize{7}{8.4}\selectfont, color=ewsslate,
      align=left] at ($(too.south east)+(8pt,-46pt)$)
  {model output ends here\\host consequences begin};

\end{tikzpicture}

\ewsdiagramnote{The campaign probes five points on one deployed pipeline rather
than five unrelated prompts. Each finding is named for the boundary it crosses,
and the dashed line marks the separation the paper holds throughout: what the
model emitted is evidence, what a host would then do is not.}
\end{ewsdiagram}

\end{ewsfigure}
\end{document}
